\chapter{La Fonction Exponentielle}

Ce chapitre consolide l'étude de la fonction exponentielle et introduit les limites de croissances comparées associées, qui constituent des outils fondamentaux pour l'étude globale des fonctions et de leurs asymptotes.

\section{Compléments sur la Fonction Exponentielle}

\subsection{Rappels et propriétés analytiques}

La fonction exponentielle, notée $\exp$ ou $x \mapsto e^x$, est l'unique fonction dérivable sur $\R$ égale à sa propre dérivée et valant 1 en 0 ($f' = f$ et $f(0)=1$).

\begin{propriete}[Propriétés fondamentales]
\begin{itemize}
    \item Pour tout réel $x$, $e^x > 0$.
    \item La fonction $x \mapsto e^x$ est strictement croissante et continue sur $\R$.
    \item Les limites aux bornes sont :
    \[
    \lim_{x \to -\infty} e^x = 0 \quad \text{et} \quad \lim_{x \to +\infty} e^x = +\infty
    \]
    Ainsi, la droite d'équation $y=0$ est une asymptote horizontale à la courbe en $-\infty$.
\end{itemize}
\end{propriete}

\subsection{Croissances comparées}

À l'infini, la fonction exponentielle croît beaucoup plus vite que n'importe quelle fonction puissance.

\begin{theoreme}[Croissances comparées pour $\exp$]
Pour tout entier naturel $n \ge 1$ :
\[
\lim_{x \to +\infty} \frac{e^x}{x^n} = +\infty \quad \text{et} \quad \lim_{x \to -\infty} x^n e^x = 0
\]
\end{theoreme}



\section{Exercices d'Entraînement Résolus}

\subsection*{Exercice 1 : Équation exponentielle simple}
Résoudre dans $\R$ l'équation suivante :
\[
e^{3x - 1} = 2
\]

\textbf{Solution :}
Puisque $2 > 0$ et la fonction $\ln$ est la bijection réciproque de la fonction exponentielle :
\[
e^{3x - 1} = 2 \iff 3x - 1 = \ln(2) \iff 3x = \ln(2) + 1 \iff x = \frac{\ln(2) + 1}{3}
\]
L'unique solution est :
\[
S = \left\{ \frac{\ln(2) + 1}{3} \right\}
\]

\subsection*{Exercice 2 : Équation du second degré en $e^x$}
Résoudre dans $\R$ l'équation suivante :
\[
e^{2x} - 4e^x + 3 = 0
\]

\textbf{Solution :}
En remarquant que $e^{2x} = (e^x)^2$, posons le changement de variable $X = e^x$. L'équation devient :
\[
X^2 - 4X + 3 = 0
\]
C'est une équation du second degré dont le discriminant est $\Delta = (-4)^2 - 4(1)(3) = 16 - 12 = 4 > 0$.
Les solutions pour $X$ sont :
\[
X_1 = \frac{4-2}{2} = 1 \quad \text{et} \quad X_2 = \frac{4+2}{2} = 3
\]
Revenons à la variable $x$ sachant que $X = e^x$ :
- $e^x = 1 \iff x = \ln(1) = 0$
- $e^x = 3 \iff x = \ln(3)$
L'ensemble des solutions de l'équation est :
\[
S = \{0; \ln(3)\}
\]

\subsection*{Exercice 3 : Limite par croissance comparée}
Déterminer la limite en $+\infty$ de la fonction $f$ définie sur $]0, +\infty[$ par :
\[
f(x) = \frac{e^x}{x^3}
\]

\textbf{Solution :}
C'est une forme indéterminée du type « $\frac{\infty}{\infty}$ ».
D'après le cours, le théorème des croissances comparées stipule que pour tout entier $n \ge 1$ :
\[
\lim_{x \to +\infty} \frac{e^x}{x^n} = +\infty
\]
En choisissant $n = 3$, on obtient directement :
\[
\lim_{x \to +\infty} f(x) = +\infty
\]

\subsection*{Exercice 4 : Limite en $-\infty$ par croissance comparée}
Déterminer la limite en $-\infty$ de la fonction $f$ définie sur $\R$ par :
\[
f(x) = x^2 e^x
\]

\textbf{Solution :}
C'est une forme indéterminée du type « $+\infty \times 0$ ».
D'après le théorème des croissances comparées, pour tout entier $n \ge 1$ :
\[
\lim_{x \to -\infty} x^n e^x = 0
\]
En posant $n = 2$, on obtient directement :
\[
\lim_{x \to -\infty} f(x) = 0
\]

\subsection*{Exercice 5 : Dérivation d'une exponentielle composée}
Déterminer la dérivée de la fonction $f$ définie sur $\R$ par $f(x) = e^{-x^2 + x}$.

\textbf{Solution :}
La fonction $f$ est de la forme $e^u$ avec $u(x) = -x^2 + x$.
La fonction $u$ est polynomiale donc dérivable sur $\R$ avec $u'(x) = -2x + 1$.
D'après la formule $(e^u)' = u' e^u$, on obtient :
\[
f'(x) = (-2x + 1)e^{-x^2 + x}
\]

\subsection*{Exercice 6 : Propriété de positivité}
Démontrer que pour tout réel $x$, la fonction $f(x) = \frac{e^x - 1}{e^x + 1}$ est bornée par $-1$ et $1$.

\textbf{Solution :}
On cherche à prouver que pour tout réel $x$ : $-1 < f(x) < 1$.
- Montrons que $f(x) < 1$ :
  \[
  f(x) - 1 = \frac{e^x - 1}{e^x + 1} - 1 = \frac{e^x - 1 - (e^x + 1)}{e^x + 1} = \frac{-2}{e^x + 1}
  \]
  Comme $e^x > 0 \implies e^x + 1 > 1 > 0$, le dénominateur est strictement positif. Le numérateur étant négatif, la différence est strictement négative, d'où $f(x) < 1$.
- Montrons que $f(x) > -1$ :
  \[
  f(x) - (-1) = \frac{e^x - 1}{e^x + 1} + 1 = \frac{e^x - 1 + e^x + 1}{e^x + 1} = \frac{2e^x}{e^x + 1}
  \]
  Le numérateur et le dénominateur sont strictement positifs, d'où la différence est strictement positive, ce qui implique $f(x) > -1$.
La fonction $f$ est donc bien bornée entre $-1$ et $1$.

\subsection*{Exercice 7 : Limite en $-\infty$ sans indétermination}
Déterminer la limite en $-\infty$ de la fonction $f$ définie sur $\R$ par :
\[
f(x) = \frac{2}{3 + e^x}
\]

\textbf{Solution :}
On sait que $\lim_{x \to -\infty} e^x = 0$.
Par limite de somme, le dénominateur tend vers $3 + 0 = 3$.
Par limite de quotient :
\[
\lim_{x \to -\infty} f(x) = \frac{2}{3}
\]

\section{Applications et Modélisation}

\subsection*{Problème 1 : Modélisation de la croissance d'une population de bactéries}
Dans un milieu de culture de laboratoire, la population de bactéries (en milliers) au cours du temps $t$ (en heures) est modélisée par la fonction $N(t) = N_0 e^{kt}$, où $N_0$ est la population initiale et $k$ une constante réelle positive.
\begin{enumerate}
    \item La population de bactéries initiale est de $1000$ bactéries (soit $N_0 = 1$). Exprimer $N(t)$ en fonction de $k$ et $t$.
    \item Au bout de 2 heures, la population a doublé et atteint $2000$ bactéries. Démontrer que la constante $k$ vaut $\frac{\ln(2)}{2}$.
    \item En déduire la population de bactéries au bout de 5 heures (arrondir à l'entier le plus proche).
    \item Au bout de combien de temps la population de bactéries dépassera-t-elle $100\,000$ bactéries ($N(t) > 100$) ? Donner la valeur exacte puis une approximation en heures et minutes.
\end{enumerate}

\textbf{Solution :}
\begin{enumerate}
    \item Avec $N_0 = 1$, la population (en milliers) est modélisée par : $N(t) = e^{kt}$.
    
    \item On sait que $N(2) = 2 \iff e^{2k} = 2$.
    En appliquant la fonction logarithme :
    \[
    2k = \ln(2) \iff k = \frac{\ln(2)}{2} \approx 0,3466
    \]
    
    \item L'expression générale de la population est : $N(t) = e^{\frac{\ln(2)}{2} t}$.
    Au bout de 5 heures ($t = 5$) :
    \[
    N(5) = e^{2,5 \ln(2)} = e^{\ln(2^{2,5})} = 2^{2,5} = \sqrt{32} \approx 5,657 \text{ milliers}
    \]
    La population de bactéries après 5 heures est d'environ 5657 bactéries.
    
    \item On cherche à résoudre $N(t) > 100 \iff e^{\frac{\ln(2)}{2} t} > 100$.
    En appliquant le logarithme :
    \[
    \frac{\ln(2)}{2} t > \ln(100) \iff t > \frac{2 \ln(100)}{\ln(2)} = \frac{4 \ln(10)}{\ln(2)}
    \]
    Calculons la valeur approchée :
    \[
    t > \frac{2 \times 4,60517}{0,69315} \approx 13,29 \text{ heures}
    \]
    Exprimons la fraction en minutes : $0,29 \text{ h} \times 60 \text{ min/h} \approx 17 \text{ minutes}$.
    La population de bactéries dépassera $100\,000$ après environ 13 heures et 17 minutes.
\end{enumerate}

\subsection*{Problème 2 : Charge d'un condensateur dans un circuit RC}
En électricité, la charge d'un condensateur de capacité $C$ à travers une résistance $R$ par un générateur de tension constante $E$ est régie par l'équation différentielle :
\[
R C \frac{\dd u_c}{\dd t} + u_c = E
\]
où $u_c(t)$ représente la tension aux bornes du condensateur au temps $t \ge 0$.
À l'instant initial $t=0$, le condensateur est déchargé, donc $u_c(0) = 0$.
\begin{enumerate}
    \item Mettre l'équation différentielle sous la forme $y' = ay + b$.
    \item Résoudre cette équation différentielle avec la condition initiale $u_c(0) = 0$.
    \item On pose $\tau = R C$ (la constante de temps du circuit). Écrire l'expression simplifiée de $u_c(t)$ en fonction de $\tau$.
    \item Calculer la limite de $u_c(t)$ lorsque $t \to +\infty$. Commenter physiquement.
    \item Calculer la valeur exacte de $u_c(\tau)$ en fonction de $E$. Quel pourcentage de la tension maximale $E$ est atteint à cet instant ?
\end{enumerate}

\textbf{Solution :}
\begin{enumerate}
    \item L'équation s'écrit $R C u_c'(t) + u_c(t) = E$. En divisant par $RC$ (qui est strictement positif) :
    \[
    u_c'(t) = -\frac{1}{R C} u_c(t) + \frac{E}{R C}
    \]
    C'est de la forme $y' = ay+b$ avec $a = -\frac{1}{R C}$ et $b = \frac{E}{R C}$.
    
    \item La solution générale de cette équation est :
    \[
    u_c(t) = K e^{-\frac{1}{R C} t} - \frac{b}{a} = K e^{-\frac{t}{R C}} - \frac{\frac{E}{R C}}{-\frac{1}{R C}} = K e^{-\frac{t}{R C}} + E
    \]
    Utilisons la condition initiale $u_c(0) = 0$ :
    \[
    u_c(0) = 0 \iff K e^{0} + E = 0 \iff K + E = 0 \iff K = -E
    \]
    La tension aux bornes du condensateur au cours du temps est donc :
    \[
    u_c(t) = E \left( 1 - e^{-\frac{t}{R C}} \right)
    \]
    
    \item En posant $\tau = R C$ :
    \[
    u_c(t) = E \left( 1 - e^{-\frac{t}{\tau}} \right)
    \]
    
    \item Puisque $\tau > 0$, $\lim_{t \to +\infty} -\frac{t}{\tau} = -\infty \implies \lim_{t \to +\infty} e^{-\frac{t}{\tau}} = 0$.
    Par suite, $\lim_{t \to +\infty} u_c(t) = E$.
    En régime permanent (lorsque le temps devient très grand), le condensateur est complètement chargé et sa tension atteint la tension du générateur $E$.
    
    \item Pour $t = \tau$ :
    \[
    u_c(\tau) = E \left( 1 - e^{-\frac{\tau}{\tau}} \right) = E(1 - e^{-1}) = E\left(1 - \frac{1}{e}\right)
    \]
    Puisque $e \approx 2,718$, $1 - e^{-1} \approx 1 - 0,368 = 0,632$, soit environ $63\%$ de la tension maximale $E$ est atteinte au bout d'une constante de temps $\tau$.
\end{enumerate}

\subsection*{Problème 3 : Température d'une tasse de café (Loi de Newton)}
Une tasse de café bien chaud est préparée et sa température initiale est de $80^\circ$C. On la place dans une pièce à température constante de $20^\circ$C. On note $f(t)$ la température du café (en $^\circ$C) au cours du temps $t$ (en minutes).
D'après la loi de refroidissement de Newton, le taux de variation de la température du café est proportionnel à la différence entre cette température et celle de l'air ambiant.
La fonction $f$ vérifie ainsi l'équation différentielle :
\[
(E) : \quad y' = -0,05(y - 20)
\]
\begin{enumerate}
    \item Mettre l'équation $(E)$ sous la forme $y' = ay + b$.
    \item Déterminer l'unique solution $f$ de $(E)$ avec la condition initiale $f(0) = 80$.
    \item Déterminer la température du café après 20 minutes (arrondir à $0,1^\circ$C près).
    \item Au bout de combien de temps la température du café descendra-t-elle en dessous de $30^\circ$C ?
\end{enumerate}

\textbf{Solution :}
\begin{enumerate}
    \item En développant l'expression de l'équation différentielle :
    \[
    y' = -0,05y + 1
    \]
    C'est de la forme $y' = ay+b$ avec $a = -0,05$ et $b = 1$.
    
    \item La solution générale de l'équation est :
    \[
    f(t) = C e^{-0,05t} - \frac{1}{-0,05} = C e^{-0,05t} + 20
    \]
    Utilisons la condition initiale $f(0) = 80$ :
    \[
    C e^0 + 20 = 80 \iff C + 20 = 80 \iff C = 60
    \]
    L'expression de la température en fonction de $t$ est donc :
    \[
    f(t) = 60e^{-0,05t} + 20
    \]
    
    \item Pour $t = 20$ minutes :
    \[
    f(20) = 60e^{-0,05 \times 20} + 20 = 60e^{-1} + 20 \approx 60 \times 0,36788 + 20 \approx 22,07 + 20 = 42,1^\circ\text{C}
    \]
    
    \item On cherche à résoudre $f(t) < 30 \iff 60e^{-0,05t} + 20 < 30$.
    \[
    60e^{-0,05t} < 10 \iff e^{-0,05t} < \frac{1}{6}
    \]
    En appliquant le logarithme (qui conserve le sens des inégalités car il est strictement croissant) :
    \[
    -0,05t < \ln\left(\frac{1}{6}\right) \iff -0,05t < -\ln(6) \iff t > \frac{\ln(6)}{0,05} = 20 \ln(6)
    \]
    Calculons la valeur numérique :
    \[
    t > 20 \times 1,79176 \approx 35,84 \text{ minutes}
    \]
    Puisque $0,84 \text{ min} \times 60 \text{ s/min} \approx 50 \text{ secondes}$, la température passera sous le seuil de $30^\circ$C après 35 minutes et 50 secondes.
\end{enumerate}

\subsection*{Problème 4 : Équation logistique de Verhulst}
Pour modéliser l'évolution d'une population de poissons dans un lac fermé, on utilise le modèle logistique de Verhulst. La population de poissons $P(t)$ (en milliers) au cours du temps $t$ (en années) vérifie l'équation différentielle :
\[
P'(t) = 0,1 P(t) \left( 1 - \frac{P(t)}{10} \right)
\]
On suppose que la population initiale est de $2000$ poissons, soit $P(0) = 2$.
\begin{enumerate}
    \item On pose la fonction $z(t) = \frac{1}{P(t)}$. Justifier que la fonction $P$ ne s'annule pas, puis montrer que la fonction $z$ vérifie l'équation différentielle linéaire :
    \[
    (E') : \quad z' = -0,1 z + 0,01
    \]
    \item Résoudre l'équation $(E')$ avec la condition initiale associée à $z(0)$.
    \item En déduire l'expression de $P(t)$ en fonction de $t$.
    \item Déterminer la limite de la population de poissons lorsque $t \to +\infty$. Commenter la signification de la valeur 10 (en milliers) dans ce modèle.
\end{enumerate}

\textbf{Solution :}
\begin{enumerate}
    \item Si $P(t_0) = 0$ pour un certain $t_0$, alors par unicité des solutions (problème de Cauchy), $P(t)$ serait identiquement nulle, ce qui contredit la condition initiale $P(0) = 2 > 0$. Ainsi, la population reste strictement positive pour tout $t \ge 0$.
    Dérivons la fonction $z(t) = \frac{1}{P(t)}$ :
    \[
    z'(t) = -\frac{P'(t)}{P(t)^2} = -\frac{0,1 P(t) \left( 1 - \frac{P(t)}{10} \right)}{P(t)^2} = - \frac{0,1 \left(1 - \frac{P(t)}{10}\right)}{P(t)} = -\frac{0,1}{P(t)} + 0,01 = -0,1 z(t) + 0,01
    \]
    La fonction $z$ est bien solution de l'équation linéaire $z' = -0,1z + 0,01$.
    
    \item L'équation différentielle $(E')$ est de la forme $z' = az+b$ avec $a = -0,1$ et $b = 0,01$.
    La solution générale de $(E')$ est :
    \[
    z(t) = C e^{-0,1t} - \frac{0,01}{-0,1} = C e^{-0,1t} + 0,1
    \]
    La condition initiale est $z(0) = \frac{1}{P(0)} = \frac{1}{2} = 0,5$.
    \[
    z(0) = 0,5 \iff C e^{0} + 0,1 = 0,5 \iff C = 0,4
    \]
    L'unique solution de $(E')$ est donc :
    \[
    z(t) = 0,4 e^{-0,1t} + 0,1
    \]
    
    \item Par définition, $P(t) = \frac{1}{z(t)}$, d'où :
    \[
    P(t) = \frac{1}{0,4 e^{-0,1t} + 0,1} = \frac{10}{4e^{-0,1t} + 1}
    \]
    
    \item Comme $\lim_{t \to +\infty} -0,1t = -\infty \implies \lim_{t \to +\infty} e^{-0,1t} = 0$.
    Par limite de quotient, on obtient :
    \[
    \lim_{t \to +\infty} P(t) = \frac{10}{0 + 1} = 10 \text{ milliers}
    \]
    La population limite de poissons dans le lac est de $10\,000$ poissons. La valeur $10$ (en milliers) correspond à la **capacité de charge** du lac, c'est-à-dire la population maximale que l'écosystème du lac peut supporter de manière durable.
\end{enumerate}

\subsection*{Problème 5 : Équation différentielle avec second membre variable}
On considère l'équation différentielle suivante sur $\R$ :
\[
(E) : \quad y' - 2y = 3x^2 - x
\]
\begin{enumerate}
    \item Déterminer les réels $a$, $b$ et $c$ tels que la fonction polynôme $P(x) = ax^2 + bx + c$ soit une solution particulière de l'équation $(E)$.
    \item Démontrer qu'une fonction $f$ est solution de $(E)$ si et seulement si la fonction $g = f - P$ est solution de l'équation homogène associée :
    \[
    (E_0) : \quad y' - 2y = 0
    \]
    \item En déduire l'ensemble des solutions de l'équation $(E)$ sur $\R$.
    \item Déterminer l'unique solution $f$ de $(E)$ vérifiant $f(0) = 1$.
\end{enumerate}

\textbf{Solution :}
\begin{enumerate}
    \item Si $P(x) = ax^2 + bx + c$ est solution de $(E)$, elle vérifie $P'(x) - 2P(x) = 3x^2 - x$.
    Calculons $P'(x) = 2ax + b$.
    En remplaçant :
    \[
    (2ax + b) - 2(ax^2 + bx + c) = 3x^2 - x \iff -2ax^2 + (2a - 2b)x + (b - 2c) = 3x^2 - x
    \]
    Par identification des coefficients de polynômes de même degré, on obtient le système :
    \[
    \begin{cases}
    -2a = 3 \\
    2a - 2b = -1 \\
    b - 2c = 0
    \end{cases} \iff 
    \begin{cases}
    a = -\frac{3}{2} \\
    -3 - 2b = -1 \iff 2b = -2 \iff b = -1 \\
    -1 - 2c = 0 \iff c = -\frac{1}{2}
    \end{cases}
    \]
    La fonction polynôme cherchée est donc :
    \[
    P(x) = -\frac{3}{2}x^2 - x - \frac{1}{2}
    \]
    
    \item - **Sens direct** : Supposons $f$ solution de $(E)$, donc $f' - 2f = 3x^2 - x$.
      Comme $P$ est également solution de $(E)$, on a $P' - 2P = 3x^2 - x$.
      En soustrayant membre à membre :
      \[
      (f' - P') - 2(f - P) = 0 \iff (f-P)' - 2(f-P) = 0
      \]
      La fonction $g = f-P$ est donc bien solution de $(E_0)$.
      - **Sens réciproque** : Supposons $g = f-P$ solution de $(E_0)$, donc $g' - 2g = 0$.
      Comme $f = g + P$, calculons :
      \[
      f' - 2f = (g' + P') - 2(g + P) = (g' - 2g) + (P' - 2P) = 0 + 3x^2 - x = 3x^2 - x
      \]
      La fonction $f$ est donc bien solution de $(E)$.
      
    \item D'après le cours, les solutions de l'équation homogène $(E_0) : y' = 2y$ sont les fonctions de la forme $g(x) = C e^{2x}$ avec $C \in \R$.
    Puisque $f = g + P$, la solution générale de l'équation $(E)$ est :
    \[
    f(x) = C e^{2x} - \frac{3}{2}x^2 - x - \frac{1}{2}
    \]
    
    \item Déterminons la constante $C$ avec la condition initiale $f(0) = 1$ :
    \[
    f(0) = 1 \iff C e^{0} - \frac{3}{2}(0)^2 - (0) - \frac{1}{2} = 1 \iff C - \frac{1}{2} = 1 \iff C = \frac{3}{2}
    \]
    L'unique solution est donc la fonction $f$ définie sur $\R$ par :
    \[
    f(x) = \frac{3}{2}e^{2x} - \frac{3}{2}x^2 - x - \frac{1}{2}
    \]
\end{enumerate}

